\runninghead{Three Readings Compared}
\section{One Kingship Read Three Ways}
\label{sec:comparison}

The three readings answer three questions that the same texts raise, and each
lets a different element lead. The first answers from the Son's humanity: how
he who is King by nature receives the kingdom as man. The second answers from
the kingdom's origin: how a kingdom not from here can be everywhere and claim
every ruler. The third answers from its end: how a kingdom handed over to the
Father can last for ever.

\begin{comparisontable}{}{The Lamb receives the kingdom}{Not from hence, and therefore everywhere}{The kingdom handed over}
Question &
How does the Lamb receive what the Son always had? &
Where is a kingdom not of this world, and how does it claim kings and nations? &
How can a kingdom handed over have no end?\\
Leading texts &
Introit, Offertory &
Gospel, Gradual, Collect &
Preface, Alleluia, Communion\\
Answer &
King by nature as God, he receives as man, through the cross, the same
dominion: one dominion by two rights &
Its origin and means are from above, so it is here in the faithful of every
nation; kings serve it freely, in their own office, and all at the judgment &
The kingdom handed over is the redeemed, brought to the vision of God; the
session begins in the cleansed soul and is completed after the judgment\\
Principal witnesses &
Hilary, Aquinas, Bede &
Augustine, Chrysostom, Bellarmine &
Augustine, Hilary\\
What it asks of the faithful &
To ask, and to share the King's self-offering &
To hear the truth and serve him in one's own office; to fight sin, not with
swords &
Purity of heart for the vision\\
Strongest objection &
Chrysostom and Cyril of Alexandria: nothing received in addition; for
Chrysostom every text of giving refers to the generation &
That the Gospel's denial leaves the subjection no public consequence, or that
the subjection makes Christ a rival of Caesar &
That the reign ends when the kingdom is handed over: the teaching Cyril of
Jerusalem opposes, and Jerome answers from Dan 7:14\\
Differences among its witnesses &
On Ps 2:7: Hilary reads the resurrection, Augustine the eternal generation;
Theodoret applies the verse to the humanity &
On Jn 18:36, the people ruled (Augustine) or the royal power (Chrysostom); on
Jn 18:37, the birth from the Virgin (Augustine) or the eternal generation
(Chrysostom) &
Hilary and Augustine on the Son's subjection; Augustine and Bellarmine on
when he sits\\
\end{comparisontable}

All three identify the kingdom with the redeemed people, the Church, and none
of them lets it become a territory. The cross has a place in each, and a
different one. In the first reading it is the means by which the kingdom is
received: the asking of the second psalm is the passion, and the Secret's
victim is the same sacrifice offered now. The second reading's centre is the
denial before Pilate and the service of kings; the cross enters it as the
standard under which the faithful fight, in Schuster's note on the
Postcommunion, and their warfare is against sin. In the third it is the
priestly offering that the Preface sets before the handing over. St Augustine
contributes to all three readings: the Psalm~2 exposition to the first, on the
decree and the asking, and to the second, on the kings; the tractate on John,
the Psalm~71 exposition and \work{Letter}~185 to the second; the \work{De
Trinitate} and the Psalm~28 exposition to the third. St Hilary also
contributes to all three: his tractate on Psalm~2 to the first, on the asking
and the inheritance, and to the second, on the kingdom where sin is overcome;
his \work{De Trinitate}, with the same tractate, to the third. St Thomas and St
Robert Bellarmine contribute to all three.

The answers differ, and each would mislead on its own. The first reading,
taken alone, would mislead if it let the kingship begin at Easter. The second,
taken alone, could leave the Preface's handing over and the Communion's
eternal session as ornaments of a kingdom wholly present in history. The
third, taken alone, could pass over the Gospel's present, mixed kingdom, with
its tares, and the service that kings owe now.

The first and second readings differ most over the Gospel, which both read of
the same prisoner. For the first, the bound King is acquiring the dominion of
the second psalm in the trial itself: he is captive, in Chrysostom's words,
``by consenting, and giving Himself up,'' and the kingship of Ps 2:8 is won by
the passion that has begun. For the second, the same prisoner declares a
kingdom whose power is from above and is ``sufficient for itself, needing
nothing,'' as Chrysostom says, a King as God ``whether man so will, or no,'' as
Cyril of Alexandria says. The first reading speaks of a dominion received, the
second of a dominion that needs nothing given. Jn 18:37 is where the two
meet. ``For this was I born'' divides the second reading's own principal
witnesses: Augustine reads it of the birth from the Virgin, Chrysostom of the
eternal generation, and Aquinas gives both. The second reading rests on
verse~36 and on the hearers of the truth in verse~37, and either birth serves
it. For the first reading, Chrysostom's answer is its strongest objection:
born a king, the Son ``has nothing which He received in addition.'' The first
reading grants that premise and refers the receiving to the humanity. Taken
with Chrysostom, then, the verse states the first reading's premise and its
objection at once; taken with Augustine, it speaks of the birth in which, on
the first reading, the kingship is received. Augustine reads Ps 2:7 of the
eternal generation and Jn 18:37 of the birth from the Virgin: the eternal
birth grounds the decree of the psalm, and the temporal birth the testimony
before Pilate.

Read together, the three readings say what the feast's Preface says in one
sentence: that the anointed Son reigns by right, reigns as the crucified
Priest, and reigns towards the Father, in whom his kingdom of truth and life,
holiness and grace, justice, love and peace has no end.
